\chapter{The Board Model}\label{ch05:chapter}

\section{The board as an evidence object}\label{ch05:sec:intro}

A \emph{board} is the candidate's unit of design: the complete, immutable,
identity-bearing description of one circuit-shaped cognitive-system model.
Everything the candidate can do --- check, simulate, trace, adapt, audit
--- it does \emph{to a board}, and everything it reports, it reports
\emph{about a board identified by digest}. This chapter walks through
every record type that a board contains, field by field, and then through
the two disciplines that make the walk worthwhile: the immutability
discipline, which guarantees that a board cannot change between the moment
it was checked and the moment its check report is relied upon; and the
identity discipline, which guarantees that any structural edit --- any
edit at all --- produces a detectably different board.

The design target is worth stating before the field lists. In an
evidence-bearing research program, the dangerous artifacts are not the
wrong ones but the \emph{ambiguous} ones: a design whose parameters might
have been edited after review, a netlist whose meaning depends on element
order, a ``safety'' annotation that might have been added by a later hand.
The board model removes ambiguity by construction. A board is a value, not
an object with a lifecycle; its identity is a deterministic function of
its complete structure; and its research boundary is rederived --- not
stored, not trusted, not inherited --- by every public operation that
touches it.

As everywhere in this monograph, the boundary comes first: a board is a
model description. Nothing in this chapter asserts that any board
describes a real cognitive system, a deployed model, or a physical
circuit. Section~\ref{ch05:sec:invariants} closes the chapter with the
formal statement.

\section{The immutability discipline}\label{ch05:sec:immutability}

Every record in the board model obeys the same five mechanical rules,
inherited from the type layer of Chapter~\ref{ch04:chapter} and worth
restating here because the board is where they compound.

\begin{ruleenv}[Frozen records]\label{ch05:rule:frozen}
Every board-model record is a frozen dataclass with slots: fields are
assigned once, at construction, and attribute assignment thereafter raises.
There are no setters, no builders that return partially initialized
values, and no mutable default fields.
\end{ruleenv}

\begin{ruleenv}[Exact types, no subclasses]\label{ch05:rule:exact}
Construction and every revalidation check the exact runtime type of the
record itself and of every typed field. A subclass of \texttt{Net},
\texttt{ComponentInstance}, or \texttt{Board} --- however
faithful-looking --- is rejected. This closes the classic loophole in
which a subclass overrides equality, hashing, or a property to lie about
its contents while passing type checks.
\end{ruleenv}

\begin{ruleenv}[Recursive freezing of open-ended containers]
\label{ch05:rule:freeze}
The two open-ended containers a component carries --- its parameter
mapping and its state-binding mapping --- are recursively frozen at
construction: mappings become immutable mapping proxies over detached
copies, sequences become tuples, sets become frozen sets, and every leaf
must be one of the canonically encodable value types
(Section~\ref{ch05:sec:parameters}). A metadata value that cannot be
frozen canonically is a construction error, not a tolerated opaque blob.
Detachment matters as much as freezing: the frozen copy shares no
structure with the caller's original, so a caller who retains and mutates
the dictionary it passed in mutates only its own copy.
\end{ruleenv}

\begin{ruleenv}[Revalidation on touch]\label{ch05:rule:revalidate}
Every record exposes a revalidation method that re-checks the full
well-formedness of the record and, transitively, of every child record.
Public operations call it on entry. Immutability makes revalidation
\emph{normally} redundant --- that is the point --- but the candidate does
not rely on the language to enforce what it can cheaply re-verify:
low-level attribute forcing, pickling round-trips, and hostile or buggy
callers are all caught by the re-check rather than trusted away.
\end{ruleenv}

\begin{ruleenv}[Safe references everywhere]\label{ch05:rule:refs}
Every identifier field in every record --- board, component, port, net,
test point, calibration, feedback loop --- must satisfy the versioned
safe-reference profile (ASCII, leading alphanumeric, at most 128
characters of alphanumerics and \texttt{.\ \_ : -}), with the single
audited exception of the \texttt{\$board} endpoint sentinel. Identifiers
appear in audit reports and evidence packets; the profile exists so that
no identifier can visually impersonate another under human review.
\end{ruleenv}

\section{Components}\label{ch05:sec:components}

\begin{definition}[Component instance]\label{ch05:def:component}
A \texttt{ComponentInstance} consists of seven fields:
\texttt{ref}, \texttt{kind}, \texttt{semantic\_profile},
\texttt{parameters}, \texttt{state\_bindings}, \texttt{safety\_critical},
and \texttt{finite\_transition}.
\end{definition}

\emph{\texttt{ref}} is the component's board-unique identifier, a safe
reference. Uniqueness is not enforced by the record --- a lone component
cannot know its siblings --- but by the design-rule checker, which fails
a board with duplicated references under its model-class diagnostics.
The division of labor is deliberate: records enforce everything locally
checkable; the DRC enforces everything relational.

\emph{\texttt{kind}} names the component's contract in the closed catalog
of Chapter~\ref{ch04:chapter} (in v0.1: \texttt{pass}, \texttt{scale},
\texttt{cjt}, \texttt{decision\_threshold}, \texttt{mux}, \texttt{fuse},
\texttt{delay}, \texttt{finite\_transition}). The record stores the name,
not the contract: resolution happens at check time against the exact
catalog whose digest is bound into the board identity, so a board checked
under one catalog cannot silently be re-interpreted under another.
An unknown kind is a semantics-class diagnostic, never a guess.

\emph{\texttt{semantic\_profile}} is the versioned name of the component's
declared meaning (for example \texttt{pcba8.cjt.reference.v0}). It exists
so that a later revision of a component's transfer semantics is a
\emph{different profile string} --- visible in every payload --- rather
than a silent behavioral change under an unchanged name. The DRC checks
that the instance's profile matches the catalog contract's profile for
its kind; a mismatch means the board was authored against a different
semantics than the engine implements, and that is an error, not a
warning.

\emph{\texttt{parameters}} is the recursively frozen mapping of
Section~\ref{ch05:sec:parameters}. The record freezes and detaches it;
the DRC later checks it against the kind's typed parameter contract ---
exact key set, value types, bounds, and choice sets.

\emph{\texttt{state\_bindings}} is a frozen mapping from binding names to
exact \texttt{CasState} literals (Chapter~\ref{ch04:chapter}). Bindings
are the board's semantic annotations: a junction may declare which CAS
states its cause, bias, and effect roles are intended to model. The
record validates that every bound value is an exact revalidated state
literal; the DRC validates the relational part --- that required binding
names for the kind are present, that no unknown binding names appear, and
that the bias-domain coherence rule holds where the kind demands it.

\emph{\texttt{safety\_critical}} is an exact boolean marking the component
as participating in a safety-relevant path. The flag has consequences
downstream --- safety-marked paths require marked test points and
calibration references, and unreachable-threshold analysis attaches to
them --- but the flag itself is a \emph{declaration}. The candidate's
standing rule applies: marking a component safety-critical creates
obligations, never credit. A board with safety markings and without the
corresponding probes and calibrations fails; a board with them merely has
its declarations checked for internal consistency. No configuration of
flags constitutes evidence that anything is safe.

\emph{\texttt{finite\_transition}} is either absent or an exact
\texttt{FiniteTransitionSemantics} record, described next. Only
components whose kind is adapter-eligible meaningfully carry one; the
TKS adapter of Chapter~9 refuses everything else.

\subsection{Finite transition semantics}\label{ch05:sec:finite}

\begin{definition}[Finite transition certificate]\label{ch05:def:finite}
A \texttt{FiniteTransitionSemantics} record binds: a versioned profile
string; a named source carrier (name plus tuple of tagged finite
elements); a named target carrier (likewise); a total mapping, given as a
tuple of integer indices into the target carrier, one per source element;
a \texttt{certificate\_sha256} over the canonical payload of all of the
above; and four booleans fixing the declared character of the transition
--- \texttt{total}, \texttt{deterministic}, \texttt{stateful},
\texttt{stochastic} --- of which the first two default true and the last
two default false.
\end{definition}

Field-level points. The carrier elements are \emph{tagged} finite values,
so a source carrier $\{0, 1\}$ and a source carrier
$\{\mathit{false}, \mathit{true}\}$ are distinct carriers with distinct
certificates (Chapter~\ref{ch04:chapter},
Rule~\ref{ch04:rule:typepun}); the mapping is by index, so the
transition table is meaningful only relative to the exact element order
serialized into the certificate; and the certificate is computed by the
same canonical-JSON-then-SHA-256 discipline as every other identity in
the candidate.

The record's docstring states its own limit, and the audit language of
this monograph adopts it verbatim in substance: the digest establishes
\emph{exact transition-table identity only}. It is not an external
review, not a safety license, and not a claim that the carrier labels
have CAS meaning. A certificate answers one question --- ``is this
exactly the table that was certified?'' --- and no other.

One method deserves special mention because it is a piece of honest test
infrastructure: the record offers a \texttt{with\_mapping} operation
that returns a \emph{deliberately uncertified} mutation --- the mapping
changes, the certificate does not --- documented as ``useful for negative
tests.'' Its existence lets the test suite verify that every downstream
consumer actually recomputes and rejects a stale certificate, rather than
trusting the stored digest. A discipline that cannot exhibit its own
forgeries cannot test its forgery detection.

\subsection{The catalog viewed from the board}\label{ch05:sec:catalog}

Because every component field ultimately answers to a catalog contract,
the closed v0.1 catalog deserves a board's-eye reference here;
Table~\ref{ch05:tab:catalog} summarizes it. The catalog is published
under the versioned profile
\texttt{cas-pcba-v8.closed-component-catalog.v1}, and its pinned digest
is itself revalidated: the catalog payload constructor re-checks, on
every computation, that the catalog container has not been substituted,
that every entry is an exact contract, and that each catalog key equals
its contract's declared kind. A substituted or diverged catalog is a
runtime error, not a silently different engine.

\begin{longtable}{@{}lp{4.2cm}p{5.3cm}l@{}}
\caption{The closed v0.1 component catalog, as consumed by boards. All
activation ports carry signal kind \textsc{activation} with unit
\textsc{activation\_0\_1}; decision ports carry \textsc{decision} with
\textsc{boolean}; finite-state ports carry \textsc{finite\_state} with
\textsc{none}. ``S'' marks stateful kinds; \texttt{finite\_transition}
is the sole TKS-adapter-eligible kind.}
\label{ch05:tab:catalog}\\
\toprule
Kind & Ports (in $\to$ out) & Parameters & Flags \\
\midrule
\endfirsthead
\toprule
Kind & Ports (in $\to$ out) & Parameters & Flags \\
\midrule
\endhead
\bottomrule
\endfoot
\texttt{pass} & \texttt{input} $\to$ \texttt{output} & (none) & \\
\texttt{scale} & \texttt{input} $\to$ \texttt{output} &
\texttt{factor} (decimal $\ge 0$); \texttt{ceiling} (decimal in
$[0,1]$) & \\
\texttt{cjt} & \texttt{cause}, \texttt{bias} $\to$ \texttt{effect} &
\texttt{beta} (decimal $\ge 0$); \texttt{saturation}, \texttt{cutoff}
(decimal in $[0,1]$); required state bindings \texttt{cause},
\texttt{bias}, \texttt{effect} & \\
\texttt{decision\_threshold} & \texttt{input} $\to$ \texttt{decision}
(boolean) & \texttt{threshold} (decimal in $[0,1]$) & \\
\texttt{mux} & \texttt{when\_false}, \texttt{when\_true},
\texttt{select} (boolean) $\to$ \texttt{output} & (none) & \\
\texttt{fuse} & \texttt{input} $\to$ \texttt{output} &
\texttt{threshold} (decimal in $[0,1]$); \texttt{action} (string in
\{\texttt{permanent\_block}, \texttt{warn}\}) & S \\
\texttt{delay} & \texttt{input} $\to$ \texttt{output} &
\texttt{initial} (decimal in $[0,1]$) & S \\
\texttt{finite\_transition} & \texttt{source} $\to$ \texttt{target}
(finite state) & (none) & TKS \\
\end{longtable}

Reading the table board-side yields several observations that recur in
the audit chapters. The junction's required state bindings are the only
place the catalog \emph{demands} CAS annotation --- a board may not
declare a \texttt{cjt} without saying which states its three roles are
intended to model, which is what makes the bias-domain coherence check
well-posed. The two stateful kinds are exactly the two whose semantics
involve memory: the unit delay (one tick of activation history, seeded
by \texttt{initial}) and the fuse (run-local open/closed state under
\texttt{permanent\_block}); statefulness is declared in the contract so
that the DRC's feedback analysis and the simulator's scheduling can
refuse configurations they do not support rather than mis-simulate
them. And the finite-transition kind is the single, narrow doorway to
the TKS adapter --- eligibility is a catalog flag, not a property a
board can assert about an arbitrary component.

\section{Typed ports}\label{ch05:sec:ports}

A board carries a tuple of \texttt{PortSpec} records
(Chapter~\ref{ch04:chapter}, Definition~\ref{ch04:def:portspec}) as its
external surface: the points at which fixtures inject inputs and observe
outputs. Each port declares its direction, signal kind, unit, and
required flag; board ports are addressed by nets through the
\texttt{\$board} endpoint sentinel, so ``the board's own port \texttt{x}''
and ``component \texttt{c}'s port \texttt{x}'' are structurally different
endpoints and can never be conflated by a name collision.

Two board-level remarks. First, the \texttt{required} flag on an input
port is the hook for the DRC's connectivity rule: a required input must
be either connected to a driving net or covered by an executable default
in the selected simulator profile; an unconnected required input with no
default is a net-class error. Second, board ports are the \emph{only}
sanctioned crossing of the board's evidential boundary. Fixtures deliver
typed, canonically serialized input samples at input ports, and every
observation in a trace is attributed either to a declared port or to a
declared test point --- there is no side channel by which a value can
enter or leave a simulation unrecorded.

\section{Parameters and finite-state annotations}
\label{ch05:sec:parameters}

The parameter and state-binding mappings are where board authors have the
most freedom, so they are where the model is most careful. Freedom of
\emph{content} is preserved: a parameter mapping may nest mappings,
sequences, and sets, and may carry decimals, integers, booleans, strings,
state literals, enumeration members, floats, and explicit nulls. Freedom
of \emph{form} is not: every leaf must be one of the canonically
encodable types, and the encoding is injective and type-tagged.

\begin{construction}[Canonical value encoding]\label{ch05:con:encoding}
The canonical encoding maps: null to a tagged none; booleans to a tagged
boolean; integers to a tagged exact decimal string; floats to a tagged
binary64 encoding by exact hexadecimal representation, with explicit
markers for \texttt{NaN} and signed infinities; decimals to a tagged
canonical string; state literals to a tagged record of all four factor
fields with explicit nulls for absent Drive and Register; enumeration
members to a tagged record of their class path and encoded value;
strings to a tagged string; mappings to a tagged, key-sorted item list
(string keys only); sequences to a tagged item list preserving order;
and sets to a tagged item list sorted by each item's canonical JSON text,
so that set encoding is independent of iteration order. Any other type
raises a board-identity error naming the offending type.
\end{construction}

Three properties of this encoding do real work. \emph{Injectivity across
types}: the tags separate integer \texttt{1} from boolean \emph{true}
from string \texttt{"1"} from decimal \texttt{Decimal("1")} from binary64
\texttt{1.0} --- five different tagged payloads, five different digests
(Chapter~\ref{ch04:chapter},
Proposition~\ref{ch04:prop:injective}). \emph{Order discipline}: mapping
items are sorted by key and set items by canonical text, so authoring
order cannot leak into identity, while sequence order --- which is
semantically meaningful --- is preserved exactly. \emph{Totality by
refusal}: the encoding is total on the declared value universe and
undefined --- loudly --- outside it, so no board can carry an
unserializable payload whose identity would then depend on ambient
object state.

State bindings ride the same machinery with one addition: every bound
\texttt{CasState} is revalidated during freezing and again during
canonical encoding, so a binding that was somehow corrupted after
construction cannot reach the identity computation intact.

The phrase \emph{finite-state annotations} in the chapter contract
deserves a precise gloss. Bindings and parameters are
\emph{annotations}: declared, structural, design-value metadata in the
sense of the measurement-kind taxonomy
(Chapter~\ref{ch04:chapter}, Section~\ref{ch04:sec:measkinds}). They
assert what the author \emph{intends} a component to model, and the DRC
checks their internal and relational coherence. They are not
observations, they carry no measurement kind other than design value,
and no accumulation of annotations upgrades itself into evidence about
any system outside the board.

\section{Nets and endpoints}\label{ch05:sec:nets}

\begin{definition}[Net]\label{ch05:def:net}
A \texttt{Net} consists of a safe reference \texttt{ref}, a tuple of
\texttt{Endpoint} records, and an exact boolean
\texttt{safety\_critical}. An \texttt{Endpoint} is an
(\texttt{owner\_ref}, \texttt{port\_ref}) pair, where
\texttt{owner\_ref} is either a component reference or the board
sentinel \texttt{\$board}.
\end{definition}

The record level enforces locality: endpoints are exact typed pairs, the
tuple is frozen, and every reference is profile-checked. Everything
relational belongs to the DRC's net-class diagnostics, and the board
model is designed so that each such rule is decidable from the frozen
structure alone: every endpoint must resolve to an existing owner and an
existing port on that owner; every net must have exactly one driving
(output-direction) endpoint; every net must have at least one load; the
driver's signal kind and unit must agree with every consumer's under the
closed admissibility table; no port may appear on two nets; and the
signal-kind/unit pair itself must be admissible. None of these checks
consults anything outside the board and the versioned catalog --- which
is what makes the check report a durable artifact rather than a
statement about a moment.

The \texttt{safety\_critical} flag on a net mirrors the component flag:
it marks the net as part of a declared safety path, obligating a marked
test point on that net (a test-point-class diagnostic if absent) and
enrolling the net's path in threshold-reachability analysis. As with
components, the flag is obligation, not credit.

\section{Test points and calibration references}\label{ch05:sec:tp}

\begin{definition}[Test point, restated at board level]
\label{ch05:def:testpoint}
A \texttt{TestPoint} binds a safe reference, the reference of the net it
observes, an optional calibration reference, and an exact boolean
\texttt{safety\_critical}.
\end{definition}

A test point is a \emph{declared observation site}: a commitment, made in
the design, that a particular net is observable and that observations of
it will be attributed to this named point. During simulation, test-point
observations enter the trace as measurement records attributed to a
producer reference derived from the test point, so the provenance of
every observed value includes \emph{where in the design} it was observed,
not merely what its value was.

The attribution vocabulary is itself typed. A measurement's producer
reference admits exactly three forms: a component reference (a safe
reference naming a declared component); the board sentinel
\texttt{\$board} (for values entering or leaving the board surface); or
a prefixed test-point form, \texttt{test-point:}\emph{ref}, whose
suffix must itself be a safe reference naming a declared test point.
Any other producer string is rejected when the measurement is
constructed or revalidated. The three forms partition the possible
origins of a value in a run --- a component computed it, the fixture
supplied it at the surface, or a probe observed it --- and the
partition is machine-checkable, so an audit can classify every number
in a trace by origin without heuristics. Nothing else can produce a
value; in particular, there is no ``anonymous'' or ``synthetic''
producer form, which is another face of the no-side-channel rule of
Section~\ref{ch05:sec:ports}.

\begin{definition}[Calibration reference, restated at board level]
\label{ch05:def:calref}
A \texttt{CalibrationRef} binds: a safe reference and a version string;
a \texttt{profile\_sha256} (exact uppercase 64-hex); a typed measurement
kind; a typed unit; instrument and protocol identity strings; a finite
decimal uncertainty; and a boolean \texttt{packet\_eligible} that
defaults to --- and in every locally constructed record remains ---
false.
\end{definition}

The field set is the candidate's answer to a specific failure mode:
calibration claims that are vague enough to be unfalsifiable. Here a
calibration reference must say \emph{which} profile (by digest),
\emph{which} instrument identity, \emph{which} protocol identity,
\emph{what} uncertainty, and \emph{what kind} of measurement it purports
to describe --- or it cannot be constructed at all. The DRC then enforces
the relational rules: an executable test point on a checked board needs a
calibration reference that exists, is well-formed, and is unit-compatible
with the observed net; safety-marked paths need their probes; and a
missing or mismatched calibration is a distinct diagnostic, never a
silent downgrade.

The \texttt{packet\_eligible} flag encodes the non-self-certification
rule at the type level: a board author can declare calibration metadata,
and the declaration is useful --- it makes the design checkable and the
fixture executable --- but nothing the author writes can mark the record
as eligible for an authenticated Evidence Packet. Eligibility is granted,
if ever, by an authentication process outside the board, against the
digest the record already carries.

\section{Declared feedback loops}\label{ch05:sec:feedback}

\begin{definition}[Feedback declaration]\label{ch05:def:feedback}
A \texttt{FeedbackLoop} binds: a safe reference; a tuple of component
references constituting the declared cycle; the reference of the delay
component that breaks the cycle; a finite decimal
\texttt{max\_round\_trip\_gain}; and a calibration reference for the
gain bound.
\end{definition}

The board model treats feedback the way it treats safety: as something
that must be \emph{declared to exist at all}. The DRC computes the
strongly connected components of the board's directed component graph;
every realizable cycle must be covered by an exact feedback declaration
naming its components; the declared delay component must be of the
stateful unit-delay kind and must lie on the cycle; the declared gain
bound must be conservative and sub-unity (with explicit declared
exceptions where the source discipline admits them); and the bound must
carry a dimensionless calibration reference. An undeclared cycle, a
cycle without delay, an unsafe gain, and a malformed declaration are four
distinct feedback-class diagnostics.

Field-level notes. The gain is a \emph{finite} decimal by type: the
idealized infinite gain that the published source uses in its safety
board (Chapter~11) is unrepresentable in a feedback declaration, exactly
as it is unrepresentable in the junction parameter contract. The
calibration reference means that even the \emph{bound} --- a design
value --- must name what kind of quantity it is and with what
uncertainty it is asserted, keeping the gain bound inside the same
measurement-kind taxonomy as every other number. And the declaration is
by component references, not by object pointers: the declared cycle is
checkable against the frozen structure, and the declaration survives
serialization unchanged.

The limits are as important as the checks, and Chapter~7 states them in
full: stateful-beyond-unit-delay loops and stochastic loops are
unsupported, and unsupported or resource-bounded analysis yields an
\emph{unresolved} safety diagnostic --- never a pass. A feedback
declaration is a claim the checker can test, not a permission slip.

\section{Source and research-boundary metadata}\label{ch05:sec:metadata}

A board's canonical payload does not begin with its components; it begins
with \emph{who is speaking}. Four blocks of metadata bind every board
identity to the exact software and arithmetic that produced and would
interpret it.

\emph{Source identity.} The versioned source-bundle digest --- a SHA-256
over the candidate package's own Python sources under the profile
\texttt{cas-pcba-v8.python-source-bundle.v1} --- enters the payload. Two
boards authored identically under different implementations have
different identities, because ``the same board'' under different
interpreting code is not evidentially the same object.

\emph{Engine identity.} The engine payload (profile
\texttt{cas-pcba-v8.reference-engine.v1}) nests the source identity, the
component-catalog digest, and the complete decimal-context payload with
its digest. The catalog digest is computed over every contract in the
closed catalog --- kinds, ports, parameter rules, bounds, choice sets,
flags --- so a catalog edit anywhere is an engine-identity change
everywhere.

\emph{Decimal context.} The full context payload (precision, rounding,
exponent range, clamp, capitals, trap list) and its digest appear
directly in the board payload as well, per
Chapter~\ref{ch04:chapter}, Rule~\ref{ch04:rule:context}: an artifact
must be interpretable without consulting any ambient arithmetic state.

\emph{Research boundary.} The fixed Ring-4 boundary travels inside the
evidence-bearing records that reference boards (measurement records,
projector identities, simulation results). The board payload itself does
not need a boundary field to be safe, because of the operational rule of
the next section --- but every artifact \emph{derived} from a board
carries the boundary explicitly, revalidated at construction.

\begin{ruleenv}[Boundary rederivation]\label{ch05:rule:boundary}
Public operations never accept a caller-supplied authority object and
never trust a stored one. Wherever a research boundary is needed, it is
derived fresh from the boundary type's fixed constructor and revalidated;
wherever one is received (inside a record), it is revalidated against the
fixed posture before use. Because the boundary type has no instance
storage (Chapter~\ref{ch04:chapter}, Definition~\ref{ch04:def:boundary}),
there is no mutable object whose corruption could propagate: the check is
against constants, every time.
\end{ruleenv}

This rule is the board model's answer to a subtle attack surface: shared
mutable ``status'' singletons. A design that stored its boundary in a
shared object would make every artifact's claim posture only as strong as
the weakest code path that could reach that object. The candidate's
design makes the posture a mathematical constant.

\section{Deterministic identity material}\label{ch05:sec:identity}

\begin{construction}[Canonical board payload]\label{ch05:con:payload}
The canonical payload of a board, under the versioned identity schema
\texttt{cas-pcba-v8.board-identity.v2}, consists of: the identity schema
name; the source identity; the engine identity; the catalog digest; the
complete decimal-context payload and digest; the board's \texttt{ref} and
\texttt{version}; its ports, sorted by reference, each with direction,
signal kind, unit, and required flag canonically encoded; its components,
sorted by reference, each as a complete canonical component payload
(itself under the versioned schema
\texttt{cas-pcba-v8.component-identity.v1}, including parameters, state
bindings, safety flag, finite-transition payload, and the decimal-context
digest); its nets, sorted by reference, with endpoints sorted by owner
then port; its test points, calibrations, and feedback loops, each sorted
by reference with all fields canonically encoded (feedback component
lists sorted for order-insensitivity of the declared set). The board
digest is the uppercase SHA-256 of the canonical JSON serialization of
this payload; a component digest is defined analogously over its own
payload.
\end{construction}

\begin{proposition}[Structural edits change the digest]
\label{ch05:prop:digest}
Let $B$ and $B'$ be boards accepted by the model. If $B$ and $B'$ differ
in any structural respect --- any field of any port, component,
parameter, state binding, net, endpoint, test point, calibration, or
feedback declaration; the board reference or version; or the interpreting
source, catalog, or decimal-context identity --- then their canonical
payloads differ, and hence their digests differ up to SHA-256 collision.
Conversely, if $B$ and $B'$ differ only in authoring order of ports,
components, nets, test points, calibrations, feedback declarations,
endpoint lists, mapping insertion order, or set iteration order, their
payloads are identical and they receive the same digest.
\end{proposition}

\begin{proof}
For the forward direction: the canonical payload embeds every structural
field through the injective tagged encoding of
Construction~\ref{ch05:con:encoding} (injectivity per
Chapter~\ref{ch04:chapter}, Proposition~\ref{ch04:prop:injective},
extended structurally to the composite encodings, whose tags separate the
container types and whose item lists preserve or canonically sort
content), so distinct structures serialize to distinct canonical texts;
SHA-256 then separates them absent collision. For the converse: every
collection in Construction~\ref{ch05:con:payload} is either sorted by a
key that the sort leaves invariant under reordering, or is a genuine
sequence whose order is itself structural (and thus out of scope of the
converse). Canonical JSON fixes key order and separators, so equal
payload objects yield byte-equal texts and equal digests.
\end{proof}

\begin{remark}[Evidence consequences]\label{ch05:rem:evidence}
Proposition~\ref{ch05:prop:digest} is the board model's contribution to
the audit trail. Every DRC report binds the digest of the exact board it
checked; every simulation result binds the exact DRC report identity and
board identity; every measurement's identity scope chains back to them.
Consequently: a post-review parameter tweak, a quietly added endpoint, a
relaxed gain bound, a re-pointed calibration --- each produces a board
whose digest matches no existing report, and every artifact that claimed
to be ``about'' the reviewed board is visibly about something else. The
converse direction is equally practical: because authoring order does not
enter identity, two independent transcriptions of the same design ---
or a serialization round trip --- agree on digest, so digest agreement is
usable as a transcription check. Both directions rest on the disclosed
external assumption of SHA-256 collision resistance
(Chapter~\ref{ch04:chapter}, Corollary~\ref{ch04:cor:digest}).
\end{remark}

\begin{remark}[What identity does not claim]
A digest identifies structure; it does not validate it. A board whose
digest is stable is not thereby checked (that is the DRC's report), let
alone correct, safe, or meaningful. Identity answers ``which board?'' ---
exactly and only that question --- so that every other question can be
asked \emph{of the right board}.
\end{remark}

\subsection{The digest chain}\label{ch05:sec:chain}

The board digest is the root of a chain that runs through every
downstream artifact. It is worth laying the chain out end to end,
because its integrity argument is compositional: each link binds the
identity of the previous link into its own canonical payload, so
tampering anywhere breaks verification everywhere downstream.

\begin{protocol}[Board-rooted evidence chain]\label{ch05:proto:chain}
\begin{enumerate}
\item \emph{Arithmetic and engine.} The decimal-context digest is a
constant of the interpreting software; the catalog digest and
source-bundle digest combine with it into the engine identity
(Section~\ref{ch05:sec:metadata}).
\item \emph{Board.} The canonical board payload embeds the engine
identity and the complete structure; the board digest identifies both
at once (Construction~\ref{ch05:con:payload}).
\item \emph{Check report.} The DRC report's identity is computed over
the exact board digest it checked, the DRC profile version, and the
complete ordered diagnostic list; ``this board passed'' is never
separable from \emph{which} board and \emph{which} checker.
\item \emph{Simulation result.} A result binds the board digest, the
exact DRC report identity, the simulator profile, the component
catalog and decimal-context digests, the source baseline, and the fixed
research boundary; its identity scope digest seeds every measurement.
\item \emph{Measurements and trace events.} Each measurement identity
is a digest over its identity scope, tick, producer, port, canonical
value text and value digest, unit, kind, calibration reference, and
child-measurement identities; trace events link measurements to ticks
under their own versioned profile.
\end{enumerate}
Verification replays the chain: recompute each canonical payload from
its stored fields, recompute each digest, and compare. A single edited
parameter at step 2 invalidates steps 3--5 wholesale; an edited
measurement value at step 5 invalidates exactly that record and every
record naming it as a child ancestor.
\end{protocol}

\begin{example}[Digest agreement as a transcription check]
\label{ch05:ex:transcription}
Illustrative: two collaborators independently transcribe the same
published reference design into board records --- one authoring
components alphabetically, the other in schematic order; one building
nets from a netlist file, the other by hand; one inserting parameters
in declaration order, the other sorted. If both transcriptions are
faithful, Proposition~\ref{ch05:prop:digest}'s converse direction says
their board digests agree, because none of those orderings enters the
canonical payload. If the digests \emph{disagree}, at least one
transcription differs structurally --- a parameter value, a missing
endpoint, a flag --- and the canonical payloads can be diffed
field-by-field to locate it. The digest thus doubles as a cheap,
deterministic cross-transcription audit, which the five-board corpus
program of Chapter~12 exploits.
\end{example}

\section{Boards as values}\label{ch05:sec:values}

It is instructive to contrast the board model with the mutable
document model that mainstream design tools --- electronic CAD suites,
diagram editors, notebook environments --- share. In the mutable model,
a design is an object with a lifecycle: it is opened, edited in place,
saved, and versioned (if at all) by timestamps or by an external
version-control system whose snapshots are taken at operator-chosen
moments. Analyses run against ``the current state,'' and the binding
between an analysis report and the design state it analyzed is
maintained by workflow discipline: the operator is trusted not to edit
between analysis and use.

The board model replaces every element of that picture. A board is a
\emph{value}: it does not open, it is constructed; it does not change,
it is superseded by a different value; it is not saved into an ambient
current state, it is referenced by digest. The binding between a report
and its subject is not workflow discipline but arithmetic
(Protocol~\ref{ch05:proto:chain}). ``Editing'' is the pure-functional
idiom: construct a new board sharing whatever substructure is unchanged
--- shared safely, because everything is immutable --- and receive a
new identity for the new whole.

Three practical consequences follow. \emph{Concurrency is free}: any
number of checks, simulations, and audits can run against the same
board value without locks, because nothing can mutate under them.
\emph{Caching is sound}: a DRC report keyed by board digest can be
reused for as long as the digest matches, with no invalidation
protocol beyond digest comparison --- and the report itself records the
checker version, so an engine upgrade naturally misses the cache.
\emph{History is cheap and honest}: an archive of board values is an
archive of every design state that ever mattered, each
self-identifying, none overwritable. The candidate's known-fault
fixtures --- boards that must fail specific checks forever
(Chapter~11) --- are simply archived values whose failing reports are
regression-tested against their unchanging digests.

The cost, equally practical, is that authoring convenience must be
rebuilt above the model: builders that accumulate records before the
single validating construction, and diff tools that compare canonical
payloads rather than object graphs. The candidate accepts the cost
deliberately. Authoring convenience above a sound core is engineering;
evidential soundness above a convenient core is, on the record of the
source audit this monograph preserves, how grammars come to not cover
their corpora and safety fuses come to not trip.

\section{Worked example fragments}\label{ch05:sec:worked}

Everything in this section is \emph{illustrative only}: fragments are
schematic, elide imports and full argument lists, and are chosen to make
the model's acceptance boundary concrete. They are not normative source
code, and no fragment asserts anything about any system outside the
candidate.

\subsection{A legal minimal board}\label{ch05:sec:legalboard}

\begin{example}[Pass-through board]\label{ch05:ex:legal}
The smallest interesting legal board wires one \texttt{pass} component
between the board surface's input and output:

\begin{itemize}
\item Board ports: \texttt{stimulus\_in} (\textsc{input},
\textsc{activation}, \textsc{activation\_0\_1}) and
\texttt{response\_out} (\textsc{output}, \textsc{activation},
\textsc{activation\_0\_1}).
\item Component: \texttt{u1} of kind \texttt{pass}, semantic profile
\texttt{pcba8.pass.v0}, empty parameters, empty state bindings.
\item Nets: \texttt{n\_in} with endpoints
(\texttt{\$board}, \texttt{stimulus\_in}) and (\texttt{u1},
\texttt{input}); \texttt{n\_out} with endpoints (\texttt{u1},
\texttt{output}) and (\texttt{\$board}, \texttt{response\_out}).
\end{itemize}

Construction succeeds; the board has a stable digest; and the DRC
accepts it: every endpoint resolves, each net has exactly one driver
(the board's input surface drives \texttt{n\_in}; \texttt{u1}'s output
drives \texttt{n\_out}), kinds and units agree, and no safety, feedback,
or calibration obligations are triggered because none are declared.
\end{example}

\begin{example}[Legal annotated junction]\label{ch05:ex:cjt}
A more representative fragment declares a junction with semantic
annotations: component \texttt{j1} of kind \texttt{cjt}, with a finite
decimal parameter \texttt{beta} of \texttt{Decimal("4")}, and state
bindings mapping \texttt{cause}, \texttt{bias}, and \texttt{effect}
binding names to exact \texttt{CasState} literals --- say
$({\wedge}8, C, \bot, \bot)$, $({\wedge}2, E, \bot, \bot)$, and
$({\wedge}9, C, \bot, \bot)$. The record freezes both mappings; the DRC
checks \texttt{beta} against the kind's nonnegative-decimal rule, checks
the binding names against the kind's required set, and applies the
bias-domain coherence rule to the bound literals. Note what has and has
not happened: the board now \emph{says} something in CAS vocabulary, and
the checker has verified that what it says is well-typed and coherent
--- and nothing has been evidenced about any actual cognitive process.
\end{example}

\begin{example}[Annotated safety path, declared correctly]
\label{ch05:ex:safety}
Illustrative: a fragment that declares a monitored refusal-shaped path
--- the shape whose published source version Chapter~11 dissects ---
with every obligation the model imposes made explicit:

\begin{itemize}
\item Components: \texttt{stage} of kind \texttt{cjt}
(\texttt{beta} $=$ \texttt{Decimal("4")}, \texttt{saturation} $=$
\texttt{Decimal("0.95")}, \texttt{cutoff} $=$ \texttt{Decimal("0.05")},
bindings on all three roles); \texttt{guard} of kind \texttt{fuse}
(\texttt{threshold} $=$ \texttt{Decimal("0.9")}, \texttt{action} $=$
\texttt{"permanent\_block"}), marked \texttt{safety\_critical}.
\item Net \texttt{n\_guard} from \texttt{stage.effect} to
\texttt{guard.input}, marked \texttt{safety\_critical}.
\item Test point \texttt{tp\_guard} on \texttt{n\_guard}, marked
\texttt{safety\_critical}, with calibration reference
\texttt{cal\_act\_v1}.
\item Calibration \texttt{cal\_act\_v1}: version, profile digest,
measurement kind \textsc{design\_value}, unit
\textsc{activation\_0\_1}, instrument and protocol identities, finite
uncertainty --- and \texttt{packet\_eligible} false, as always for a
locally declared record.
\end{itemize}

The board is record-legal, and --- with the remaining wiring in place
--- passes the test-point and calibration classes of the DRC. Whether
it passes \emph{threshold reachability} depends on the numbers: with
\texttt{saturation} $=$ \texttt{0.95} above \texttt{threshold} $=$
\texttt{0.9}, the conservative bound propagation can resolve the
crossing as reachable; had \texttt{saturation} been
\texttt{Decimal("0.1")} --- the published source's own configuration
--- the checker must flag the unreachable threshold. The example is the
board model doing its job: the difference between a checkable safety
declaration and an uncheckable safety \emph{story} is exactly the
difference between these typed fields and prose. No outcome of the
check, in either direction, is evidence that any real refusal
mechanism exists or works; it is evidence about the declared design.
\end{example}

\subsection{Reading a canonical payload}\label{ch05:sec:reading}

Auditors meet boards mostly as canonical payloads, so a reading guide
belongs in this chapter. A canonical board payload is a JSON object
with sorted keys and fixed separators; its top level always contains,
in the model's v2 schema: \texttt{identity\_schema} (the constant
\texttt{cas-pcba-v8.board-identity.v2}); \texttt{source} and
\texttt{engine} (the interpreting-software identities of
Section~\ref{ch05:sec:metadata}); \texttt{catalog\_sha256};
\texttt{decimal\_context} and \texttt{decimal\_context\_sha256};
\texttt{ref} and \texttt{version}; and the six sorted collections
(\texttt{ports}, \texttt{components}, \texttt{nets},
\texttt{test\_points}, \texttt{calibrations},
\texttt{feedback\_loops}).

Practical audit heuristics follow directly from the schema. To confirm
two payloads describe the same design under different engines, strip
the \texttt{source}, \texttt{engine}, \texttt{catalog\_sha256}, and
decimal-context members and diff the remainder --- the model
deliberately makes ``same structure, different interpreter'' a
well-defined, mechanically checkable relation, while the full digests
remain distinct (Proposition~\ref{ch05:prop:digest}). To locate a
structural difference flagged by digest mismatch, diff the sorted
collections pairwise by \texttt{ref}; because every element is
reference-sorted and canonically encoded, the first differing byte
lies inside the first genuinely differing field. And to audit a
parameter's exact type --- the difference between \texttt{1},
\texttt{true}, \texttt{"1"}, and \texttt{Decimal("1")} --- read its
tag: the payload preserves the type tags of
Construction~\ref{ch05:con:encoding} verbatim, so type puns that would
be invisible in a pretty-printed value listing are explicit in the
canonical form.

\subsection{Illegal constructions at the record layer}
\label{ch05:sec:illegalrecord}

\begin{example}[Rejected at construction]\label{ch05:ex:construct}
Each fragment below raises a typed error before any board exists:

\begin{itemize}
\item \emph{Subclassed record.} \texttt{class SneakyNet(Net): ...} ---
any instance is rejected by the exact-type check
(Rule~\ref{ch05:rule:exact}), regardless of field content.
\item \emph{Unsafe reference.} A component with
\texttt{ref="u1\textbackslash u202e"} (a bidirectional control
character) or \texttt{ref="\_hidden"} (leading non-alphanumeric) fails
the safe-reference profile.
\item \emph{Untyped gain.} A feedback declaration with
\texttt{max\_round\_trip\_gain=0.8} (binary64 float) is rejected: the
field requires an exact finite \texttt{Decimal}. So is
\texttt{Decimal("Infinity")} --- finiteness is checked, which is
precisely what makes the source's idealized infinite gain
unrepresentable.
\item \emph{Unfreezable metadata.} A parameter mapping containing a live
object --- a file handle, a lambda, an arbitrary class instance --- is
rejected by the canonical freeze
(Rule~\ref{ch05:rule:freeze}): the value has no canonical encoding, so
it may not enter a board.
\item \emph{Type-punned parameter.} A \texttt{fuse} authored with
\texttt{action=True} fails the string-typed choice rule; authored with
\texttt{action="permanent\_block"} it passes; authored with
\texttt{action="Permanent\_Block"} it fails the closed choice set ---
case is content.
\item \emph{Mutation attempt.} Given a constructed net \texttt{n},
the assignment \texttt{n.safety\_critical = True} raises: frozen
records have no assignment path
(Rule~\ref{ch05:rule:frozen}). The only way to ``edit'' a board is to
construct a new one --- which, by
Proposition~\ref{ch05:prop:digest}, has a new identity.
\end{itemize}
\end{example}

\subsection{Illegal boards at the check layer}\label{ch05:sec:illegaldrc}

\begin{example}[Type-legal, check-illegal]\label{ch05:ex:drc}
The record layer accepts each of the following boards --- every record is
locally well-formed --- and the DRC then fails each with a distinct
diagnostic, illustrating the division of labor of
Section~\ref{ch05:sec:components}:

\begin{itemize}
\item Two components both named \texttt{u1}: duplicate-reference, a
model-class error.
\item A component of kind \texttt{transistor} (not in the closed
catalog): unknown kind, a semantics-class error --- not an attempt at
nearest-match resolution.
\item A net whose endpoint names port \texttt{inptu} on \texttt{u1}:
unknown port, a net-class error; typos are errors, not warnings.
\item A net driven by two outputs: invalid driver count. A net driven by
none, or loading no input: likewise distinct net-class errors.
\item A decision output wired to an activation input: signal-kind
mismatch under the closed admissibility table; no coercion from boolean
to $[0,1]$ exists anywhere in the candidate.
\item A safety-critical net with no marked test point, or a test point
whose calibration reference names a nonexistent calibration: distinct
test-point-class and calibration-class errors.
\item A cycle of components with no feedback declaration: undeclared
feedback. With a declaration but no delay component on the cycle:
feedback without delay. With a declared gain of \texttt{Decimal("1.3")}:
unsafe feedback gain.
\end{itemize}

In every case the board \emph{exists} and has a digest --- one can refer
to it, archive it, and regression-test against it --- but it is failed,
and a failed report is bound to that digest permanently. The candidate's
known-fault fixtures (Chapter~11) rely on exactly this: a board that
must fail, kept, with its failure guarded by tests.
\end{example}

\section{Invariants and boundary}\label{ch05:sec:invariants}

\begin{invariant}[Board-model invariants]\label{ch05:inv:summary}
For every board accepted by the model:
(i) the board and every contained record are frozen, exact-typed, and
revalidatable, and all open-ended metadata is recursively frozen and
detached;
(ii) every identifier satisfies the safe-reference profile, with
\texttt{\$board} admitted only as an endpoint owner;
(iii) every numeric field that enters semantics is a finite
\texttt{Decimal}, and no non-finite value is representable in any
parameter, gain, uncertainty, or weight;
(iv) the canonical payload embeds the complete structure together with
source, engine, catalog, and decimal-context identities, and the board
digest is a deterministic function of exactly that payload;
(v) authoring order never affects identity; structural content always
does;
(vi) calibration records constructed locally are never packet-eligible,
and no board-resident flag can grant evidentiary authority; and
(vii) every public operation rederives and revalidates the fixed
research boundary rather than trusting any stored or supplied authority
object.
\end{invariant}

\begin{threat}[Residual limits]\label{ch05:threat:limits}
The invariants above are enforced by tested, not proved, code; the
bounded-evidence posture of Chapter~12 applies. Identity rests on the
disclosed SHA-256 collision-resistance assumption. Revalidation defends
against corruption it can observe at the record surface; it cannot
detect a hostile interpreter or a modified source bundle except insofar
as the source digest exposes them. And nothing in the model prevents an
author from writing a board that is well-formed, checkable, and
meaningless --- well-formedness is the floor of the discipline, not its
ceiling.
\end{threat}

\begin{boundary}\label{ch05:bnd:main}
This chapter specifies the board model of the CAS PCBA~8 v0.1 research
candidate, at Ring~4, \textsc{not licensed}. A board is a typed model
description; its digests identify software structures only. No board,
binding, flag, calibration reference, feedback declaration, or worked
example in this chapter constitutes evidence about any cognitive system,
deployed model, or physical circuit; no safety, causal, physical,
clinical, regulatory, or product claim is made or licensed; and all
example fragments are illustrative, non-normative, and assert nothing
beyond the acceptance behavior of the declared software contracts.
\end{boundary}
